%% ----------------------------------------------------------------------------
%%  §1.4  Numeri naturali, interi, razionali
%% ----------------------------------------------------------------------------
\section{Numeri naturali, interi, razionali}
\label{sec:nzq}

Abbiamo visto come tra gli assiomi dei numeri reali ci sia l'esistenza degli elementi
neutri $0$ e $1$. Quindi apparterranno ad $\R$ anche i risultati delle operazioni
eseguite a partire da $0$ e $1$: in particolare sono numeri reali $1+1=2$, $(1+1)+1=3$, e
così via.

\begin{definizione}[Numeri naturali, interi, razionali]\index{numeri naturali}\index{numeri interi}\index{numeri razionali}
Il sottoinsieme di $\R$ ottenuto a partire da $1$ sommando ripetutamente $1$ si chiama
insieme dei \emph{numeri naturali} e si indica con
\begin{equation}\label{eq:N}
  \N=\{1,2,3,\dots,n,\dots\}.
\end{equation}
Indichiamo con $\Z$ il sottoinsieme di $\R$ costituito dagli elementi di $\N$, dai loro
opposti e dallo zero; i suoi elementi si chiamano \emph{numeri interi} (o
\emph{interi relativi}):
\begin{equation}\label{eq:Z}
  \Z=\{0,\pm1,\pm2,\dots\}=\{0\}\cup\{\pm n : n\in\N\}.
\end{equation}
I risultati della divisione $m/n$ (che, con la terminologia introdotta dagli assiomi,
significa $m\cdot n^{-1}$) con $m,n\in\Z$ e $n\neq0$ si chiamano \emph{numeri razionali}
e si indicano con
\begin{equation}\label{eq:Q}
  \Q=\set{\frac{m}{n} \tc m,n\in\Z,\ n\neq0}.
\end{equation}
\end{definizione}

La frase «sommando ripetutamente $1$» è meno innocente di quanto sembri: definire $n$
come $1+1+\dots+1$, con $n$ addendi, usa proprio il numero $n$ che si vuole definire. La
renderemo rigorosa nel \cref{sec:induzione}, dove $\N$ sarà caratterizzato come il più
piccolo sottoinsieme di $\R$ che contiene $1$ e che, insieme a ogni suo elemento $n$,
contiene anche $n+1$.

Risulta $\N\subseteq\Z\subseteq\Q\subseteq\R$. Naturalmente, essendo $\N$, $\Z$, $\Q$
sottoinsiemi di $\R$, su di essi sono definite le operazioni di addizione e di
moltiplicazione e l'ordinamento indotti da $\R$. Però \emph{essi non soddisfano tutti
gli assiomi dei numeri reali}. Ad esempio $\N$ non soddisfa \ref{ass:opposto}:
nell'ambito dei soli numeri naturali non esiste l'opposto di alcun numero. $\Z$ non
soddisfa \ref{ass:inverso}: tutti i numeri interi non nulli, esclusi $\pm1$, hanno per
inverso un numero reale che non è intero.

Si può verificare che invece $\Q$ soddisfa tutte le proprietà algebriche relative alle
operazioni e all'ordine, cioè gli assiomi (A1)--(A6) e (O1)--(O4): la somma, il
prodotto, l'opposto e l'inverso di numeri razionali sono razionali. \emph{L'unico assioma
non soddisfatto da $\Q$ è l'assioma di completezza}. Questo è il risultato centrale del
paragrafo, e per dimostrarlo premettiamo un teorema celebre, che risale alla scuola
pitagorica.

\begin{teorema}[Irrazionalità di $\sqrt2$]\label{teo:radice2}\index{irrazionalità di $\sqrt2$}
Non esiste alcun numero razionale $c$ tale che $c^{2}=2$.
\end{teorema}

\begin{schema}
\Ipotesi{Nessuna, oltre agli assiomi (e alle proprietà elementari dei numeri interi,
che diamo per note: un intero è pari oppure dispari, il quadrato di un dispari è
dispari, ogni frazione si può ridurre ai minimi termini).}
\Tesi{Per ogni $c\in\Q$ risulta $c^{2}\neq2$.}
\end{schema}

\begin{strategia}
È una tesi \emph{negativa} («non esiste»), e le tesi negative si dimostrano quasi
sempre \emph{per assurdo}: supponiamo che un tale $c$ esista e facciamo discendere da
questa supposizione una contraddizione. La contraddizione sarà di tipo aritmetico:
scriveremo $c=m/n$ con la frazione ridotta ai minimi termini, così che $m$ e $n$ non
siano entrambi pari, e dalla relazione $m^{2}=2n^{2}$ dedurremo che $m$ è pari; ma
allora $m^{2}$ è divisibile per $4$, dunque $n^{2}$ è pari, dunque anche $n$ è pari.
Entrambi pari: ecco l'assurdo. L'intera dimostrazione ruota attorno a un'osservazione
semplice: \emph{il quadrato di un numero dispari è dispari}, e quindi se $m^{2}$ è pari,
$m$ è pari.
\end{strategia}

\begin{dimostrazione}
Sia per assurdo $c$ un numero razionale tale che $c^{2}=2$; possiamo supporre $c>0$
(altrimenti consideriamo $-c$, che ha lo stesso quadrato). In base alla \eqref{eq:Q}
esistono $m,n$ numeri interi, che possiamo supporre entrambi positivi, tali che
$c=m/n$. Se necessario, possiamo «semplificare» la frazione $m/n$, ottenendo $m$ e $n$
\emph{non entrambi pari}. Risulta $(m/n)^{2}=c^{2}=2$, cioè
\[
  2n^{2}=m^{2}.
\]
Essendo il primo membro $2n^{2}$ un numero intero pari, anche $m^{2}$ deve essere pari;
ma allora anche $m$ deve essere pari (se $m$ fosse dispari, $m=2h+1$, anche
$m^{2}=4h^{2}+4h+1=2(2h^{2}+2h)+1$ sarebbe dispari). Quindi $m=2k$, con $k$ intero.
Ne segue che
\begin{equation}\label{eq:radice2-passo}
  2n^{2}=m^{2}=4k^{2}, \qquad\text{cioè}\qquad n^{2}=2k^{2}.
\end{equation}
Ripetendo lo stesso ragionamento sulla \eqref{eq:radice2-passo}, risulta che anche
$n$ deve essere un numero pari, ciò che contrasta con l'ipotesi che $m$ ed $n$ siano
numeri interi non entrambi pari. L'assurdo prova la tesi.
\end{dimostrazione}

\begin{notastorica}[Ippaso e la crisi degli incommensurabili]
Il \cref{teo:radice2} è, nella sostanza, uno dei più antichi teoremi della storia
della matematica. I pitagorici (VI--V secolo a.C.) credevano che «tutto è numero», nel
senso che il rapporto tra due grandezze omogenee qualsiasi --- due segmenti, per esempio
--- fosse sempre esprimibile come rapporto di numeri interi. La scoperta che la diagonale
e il lato di un quadrato non hanno una misura comune (sono \emph{incommensurabili}: il
loro rapporto è $\sqrt2$) fu un trauma per la scuola. La tradizione, riportata da
Giamblico, attribuisce la scoperta a Ippaso di Metaponto, e racconta che i confratelli lo
avrebbero fatto annegare, o quantomeno espulso, per aver divulgato il segreto. La
dimostrazione «pari e dispari» che abbiamo dato è citata da Aristotele negli
\emph{Analitici primi} come esempio tipico di ragionamento per assurdo, ed è riportata in
appendice al libro X degli \emph{Elementi} di Euclide.
La crisi ebbe conseguenze durature. Se gli interi e le frazioni non bastano a misurare
tutte le grandezze, si può rinunciare a misurarle e confrontarne direttamente i
rapporti: è la \emph{teoria delle proporzioni}, che la tradizione attribuisce a Eudosso
di Cnido e che occupa il libro~V degli \emph{Elementi}. La definizione~5 di quel libro
dice, in linguaggio moderno, che $a:b=c:d$ quando, per ogni coppia di naturali $m,n$,
il multiplo $ma$ supera, uguaglia o è minore di $nb$ esattamente quando $mc$ supera,
uguaglia o è minore di $nd$. Poiché $ma>nb$ equivale a $\frac ab>\frac nm$, un rapporto è
individuato dal modo in cui \emph{divide} le frazioni in quelle che gli stanno sotto e
quelle che gli stanno sopra: è l'idea delle sezioni di Dedekind
(\S\ref{subsec:assioma-dedekind}), con oltre due millenni di anticipo. Mancava soltanto
l'ipotesi che \emph{ogni} taglio delle frazioni provenga da un rapporto, cioè la
completezza. Il prezzo fu una lunga separazione tra il mondo discreto dei numeri e quello
continuo delle grandezze geometriche, che si ricompose solo in età moderna.
\end{notastorica}

\begin{teorema}[$\Q$ non è completo]\label{teo:Q-non-completo}\index{completezza!di $\Q$ (mancanza)}
L'insieme $\Q$ dei numeri razionali non soddisfa l'assioma di completezza.
\end{teorema}

\begin{schema}
\Ipotesi{Gli assiomi dei numeri reali e il \cref{teo:radice2}.}
\Tesi{Esistono due sottoinsiemi non vuoti $A,B\subseteq\Q$ con $a\le b$ per ogni
$a\in A$, $b\in B$, tali che \emph{nessun} numero razionale $c$ soddisfa $a\le c\le b$
per ogni $a\in A$ e $b\in B$.}
\end{schema}

\begin{strategia}
Dobbiamo esibire un controesempio: due «mezze rette razionali» separate da un buco.
Il buco, naturalmente, sarà $\sqrt2$. Prendiamo per $A$ i razionali «a sinistra di
$\sqrt2$» e per $B$ quelli «a destra», ma --- attenzione --- senza usare $\sqrt2$, che
in $\Q$ non esiste: diremo che $a\in A$ se $a\le0$ oppure $a^{2}<2$, e che $b\in B$ se
$b>0$ e $b^{2}>2$. Per il \cref{teo:radice2} ogni razionale sta in $A$ o in $B$, e in
uno solo dei due. Se esistesse un separatore razionale $c$, dovrebbe stare in $A$
oppure in $B$. Se $c\in A$, $c$ sarebbe il massimo di $A$; ma mostreremo che $A$ non
ha massimo: per ogni $c\in A$ con $c>0$ riusciamo a trovare un razionale \emph{un po'
più grande}, $c+1/n$, che sta ancora in $A$, cioè con $(c+1/n)^{2}<2$. Il lavoro
tecnico consiste nel trovare un $n$ abbastanza grande, e lo faremo con una stima
esplicita. Il caso $c\in B$ è analogo.
\end{strategia}

\begin{dimostrazione}
Consideriamo gli insiemi
\begin{align}
  A&=\{a\in\Q : a\le0\}\cup\{a\in\Q : a>0,\ a^{2}<2\}, \label{eq:AB-radice2}\\
  B&=\{b\in\Q : b>0,\ b^{2}>2\}. \notag
\end{align}
Entrambi sono non vuoti ($0\in A$, $2\in B$). Tutti i numeri di $A$ sono minori di
tutti i numeri di $B$: se $a\le0<b$ è ovvio; se $0<a$ e fosse $a\ge b$, per la
\cref{prop:ordine-prodotto} sarebbe $a^{2}\ge ab\ge b^{2}>2$, contro $a^{2}<2$. Inoltre,
per il \cref{teo:radice2}, risulta $A\cup B=\Q$ e $A\cap B=\emptyset$: un razionale
positivo ha quadrato minore di $2$ oppure maggiore di $2$, mai uguale.

Se esistesse un numero \emph{razionale} $c$ con la proprietà che $a\le c\le b$ per ogni
$a\in A$ e $b\in B$, tale numero dovrebbe appartenere ad $A$ oppure a $B$.

\emph{Supponiamo $c\in A$.} Non potendo essere $c\le0$ (perché $c\ge1\in A$), ne segue
che $c>0$ e $c^{2}<2$. Sia $n$ un numero naturale maggiore di
$(2c+1)/(2-c^{2})$;\footnote{L'esistenza di un numero naturale maggiore di un
qualsiasi numero reale assegnato è la \emph{proprietà di Archimede}, che dimostreremo
nel prossimo capitolo come conseguenza dell'assioma di completezza (\cref{teo:archimede}). Qui, trattandosi
del numero razionale $(2c+1)/(2-c^{2})=p/q$ con $p,q\in\N$, basta prendere
$n=p+1$: infatti $p/q\le p<p+1$.} tale $n$ esiste. Allora, essendo $1/n^{2}\le1/n$,
\begin{equation}\label{eq:stima-radice2}
  \Bigl(c+\frac1n\Bigr)^{2} = c^{2}+\frac{2c}{n}+\frac{1}{n^{2}}
  \le c^{2}+\frac{2c+1}{n} < c^{2}+(2-c^{2}) = 2,
\end{equation}
dove nell'ultima disuguaglianza abbiamo usato $n>(2c+1)/(2-c^{2})$, cioè
$(2c+1)/n<2-c^{2}$. Dunque $c+\frac1n\in A$, il che è assurdo perché $c$ è maggiore o
uguale di tutti gli elementi di $A$, mentre $c+\frac1n>c$.

\emph{Supponiamo $c\in B$.} Analogamente si perviene a un assurdo: $c>0$, $c^{2}>2$, e
scelto $n$ naturale maggiore di $2c/(c^{2}-2)$ si ha
\[
  \Bigl(c-\frac1n\Bigr)^{2} = c^{2}-\frac{2c}{n}+\frac{1}{n^{2}} > c^{2}-\frac{2c}{n}
  > c^{2}-(c^{2}-2) = 2,
\]
e inoltre $c-\frac1n>0$ (perché $n>2c/(c^{2}-2)>1/c$, dato che $2c^{2}>c^{2}-2>0$);
quindi $c-\frac1n\in B$, assurdo perché $c$ è minore o uguale di tutti gli elementi di
$B$.

In conclusione nessun numero razionale separa $A$ e $B$: l'assioma di completezza non
vale in $\Q$.
\end{dimostrazione}

\begin{figure}[htb]
\centering
\begin{tikzpicture}[scale=0.93]
  \draw[-Latex, grigio] (-0.5,0) -- (11.5,0);
  % razionali in A (sotto) e in B (sopra) che si avvicinano a sqrt2
  \pgfmathsetmacro{\s}{sqrt(2)}
  \foreach \x in {0,0.25,...,1.3} \fill[blu] (\x*5,0) circle (1.2pt);
  \foreach \x in {1.5,1.75,...,2.3} \fill[rame] (\x*5,0) circle (1.2pt);
  \foreach \p/\q in {1/1, 3/2, 7/5, 17/12, 41/29, 99/70} {
    \pgfmathsetmacro{\v}{\p/\q}
    \ifdim \v pt < \s pt \fill[blu] (\v*5,0) circle (1.6pt); \else \fill[rame] (\v*5,0) circle (1.6pt); \fi}
  \foreach \x/\l in {0/0, 5/1, 10/2} \draw[grigio] (\x,0.07) -- (\x,-0.07);
  \node[below, font=\scriptsize, grigio] at (0,-0.08) {$0$};
  \node[blu, below, font=\scriptsize] at (5,-0.08) {$1$};
  \node[below, font=\scriptsize, grigio] at (10,-0.08) {$2$};
  % etichette con linee di richiamo, per non sovrapporle vicino a sqrt2
  \draw[grigio, very thin] (7.0,-0.06) -- (6.3,-0.42)
    node[blu, left, font=\scriptsize, inner sep=1pt] {$\tfrac75$};
  \draw[grigio, very thin] (7.083,0.06) -- (6.3,0.42)
    node[rame, left, font=\scriptsize, inner sep=1pt] {$\tfrac{17}{12}$};
  \node[rame, above, font=\scriptsize] at (7.5,0.08) {$\tfrac32$};
  \draw[grigio, very thin] (\s*5,-0.08) -- (8.3,-0.42)
    node[rosso, right, font=\scriptsize, inner sep=1pt] {$\sqrt2\notin\Q$};
  \draw[rosso, thick] (\s*5,0) circle (2.4pt);
  \draw[blu, thick, decorate, decoration={brace, amplitude=4pt, mirror}] (-0.3,-0.85) -- (6.95,-0.85)
    node[midway, below=5pt, font=\scriptsize] {$A=\{a\le0\}\cup\{a>0,\ a^{2}<2\}$};
  \draw[rame, thick, decorate, decoration={brace, amplitude=4pt}] (7.2,0.75) -- (11.3,0.75)
    node[midway, above=5pt, font=\scriptsize] {$B=\{b>0,\ b^{2}>2\}$};
\end{tikzpicture}
\caption{Il «buco» di $\Q$: gli insiemi $A$ e $B$ della \eqref{eq:AB-radice2} esauriscono
$\Q$, ma nessun razionale li separa. Le frazioni $1,\tfrac75,\tfrac{41}{29}\in A$ e
$\tfrac32,\tfrac{17}{12},\tfrac{99}{70}\in B$ si avvicinano quanto si vuole al punto
$\sqrt2$, che però non appartiene a $\Q$.}
\label{fig:buco-Q}
\end{figure}

\begin{osservazione}[Dove vive $\sqrt2$]
Osserviamo ora che i due insiemi $A$ e $B$, costituiti da numeri razionali, possono
essere riguardati come insiemi di numeri \emph{reali}: e quindi, per l'assioma di
completezza, esiste un numero reale $c$ con la proprietà che $a\le c\le b$ per ogni
$a\in A$, $b\in B$. Tale numero --- che si può dimostrare essere unico --- è tale che
$c^{2}=2$, è irrazionale e si denota con il simbolo $\sqrt2$. In altre parole:
\emph{l'assioma di completezza è esattamente ciò che garantisce l'esistenza di
$\sqrt2$}, e più in generale, come vedremo, l'esistenza delle radici, dei logaritmi, dei
limiti. Senza di esso i numeri reali sarebbero indistinguibili dai razionali, e la retta
sarebbe piena di buchi.
\end{osservazione}

\begin{idea}[Un riassunto in parole semplici]
Nell'insieme dei numeri naturali $\N$ si possono eseguire le operazioni di addizione e di
moltiplicazione, ma non è possibile in generale eseguire le operazioni inverse di
sottrazione e di divisione. $\Z$ è un ampliamento di $\N$ che permette di calcolare
anche le differenze, ma non i quozienti. $\Q$ è un ulteriore ampliamento: in $\Q$ è
possibile eseguire le quattro operazioni fondamentali (tranne, naturalmente, la
divisione per zero), ma non è possibile in generale eseguire altri calcoli altrettanto
utili, come ad esempio l'estrazione di radice. Come vedremo, $\R$ è invece
sufficientemente ricco per la maggior parte delle applicazioni: vi si possono
estrarre radici, calcolare logaritmi, e soprattutto \emph{passare al limite}.
\end{idea}

\begin{python}[Approssimare $\sqrt2$ con razionali, da sotto e da sopra]
# Costruiamo razionali in A (quadrato < 2) e in B (quadrato > 2)
# sempre più vicini tra loro; con Fraction restiamo sempre in Q.
from fractions import Fraction as F

a, b = F(1), F(2)              # a in A (1^2 < 2), b in B (2^2 > 2)
for passo in range(1, 31):
    m = (a + b) / 2            # punto medio: è ancora razionale
    if m * m < 2:
        a = m                  # m sta in A: alziamo a
    else:
        b = m                  # m sta in B: abbassiamo b
    if passo in (1, 2, 3, 10, 20, 30):
        print(f"passo {passo:2d}: {float(a):.12f} < √2 <",
              f"{float(b):.12f}   (b - a = {float(b - a):.2e})")
# La distanza b - a si dimezza a ogni passo e tende a zero, ma
# nessun razionale ha quadrato uguale a 2 (irrazionalità di √2):
print(any(x * x == 2 for x in (a, b)))
\end{python}
\begin{risultato}
passo  1: 1.000000000000 < √2 < 1.500000000000   (b - a = 5.00e-01)
passo  2: 1.250000000000 < √2 < 1.500000000000   (b - a = 2.50e-01)
passo  3: 1.375000000000 < √2 < 1.500000000000   (b - a = 1.25e-01)
passo 10: 1.414062500000 < √2 < 1.415039062500   (b - a = 9.77e-04)
passo 20: 1.414213180542 < √2 < 1.414214134216   (b - a = 9.54e-07)
passo 30: 1.414213561453 < √2 < 1.414213562384   (b - a = 9.31e-10)
False
\end{risultato}

Il programma illustra il contenuto del \cref{teo:Q-non-completo}: le due successioni di
razionali si «stringono» attorno a un punto, ma quel punto non è razionale. Il metodo
usato (dimezzare ripetutamente un intervallo) si chiama \emph{bisezione} e lo
ritroveremo nel capitolo sulle funzioni continue, dove sarà alla base della
dimostrazione del teorema dell'esistenza degli zeri (\cref{teo:zeri}).

\begin{esercizio}
Dimostrare che $\sqrt3$ è irrazionale, adattando la dimostrazione del
\cref{teo:radice2}. (\emph{Suggerimento:} al posto di «pari/dispari» si usi la
divisibilità per $3$: se $m^{2}$ è divisibile per $3$ allora $m$ è divisibile per $3$,
perché se $m=3h\pm1$ allora $m^{2}=3(3h^{2}\pm2h)+1$.)
\begin{soluzione}
Supponiamo per assurdo che esista $c\in\Q$, $c>0$, con $c^{2}=3$, e scriviamo
$c=m/n$ con $m,n\in\N$ non entrambi divisibili per $3$ (basta semplificare la
frazione). Allora $m^{2}=3n^{2}$, e $m^{2}$ è divisibile per $3$. Ogni intero è della
forma $3h$, $3h+1$ oppure $3h-1$; negli ultimi due casi
$m^{2}=9h^{2}\pm6h+1=3(3h^{2}\pm2h)+1$ non è divisibile per $3$. Dunque $m=3k$ con
$k\in\N$, e da $3n^{2}=9k^{2}$ segue $n^{2}=3k^{2}$: per lo stesso motivo anche $n$ è
divisibile per $3$, contro la scelta di $m$ e $n$. \emph{Scelta logica:} la
dimostrazione del \cref{teo:radice2} usa soltanto la proprietà «se $2$ divide $m^{2}$,
allora $2$ divide $m$»; sostituendo la classificazione pari/dispari con quella dei resti
della divisione per $3$, il ragionamento per assurdo si ripete parola per parola.
\end{soluzione}
\end{esercizio}

\begin{esercizio}
Dimostrare che la somma di un numero razionale e di un numero irrazionale è irrazionale.
Che cosa si può dire della somma di due irrazionali?
\begin{soluzione}
Sia $q\in\Q$ e $x\notin\Q$. Se $q+x$ fosse razionale, lo sarebbe anche
$(q+x)-q=x$, perché $\Q$ è chiuso rispetto alla differenza: assurdo. La somma di due
irrazionali può invece essere razionale: $\sqrt2+(-\sqrt2)=0$. \emph{Scelta logica:}
«irrazionale» è una proprietà negativa («non razionale»), che si dimostra per assurdo,
come nel \cref{teo:radice2}; la chiusura di $\Q$ rispetto alla differenza trasforma
l'ipotesi assurda in una contraddizione immediata. Per la seconda domanda, che chiede di
smentire una possibile regola generale, basta un esempio.
\end{soluzione}
\end{esercizio}

\subsection{Approfondimento: l'assioma di Dedekind}
\label{subsec:assioma-dedekind}

L'assioma di completezza si trova enunciato, nei vari testi, in forme diverse. Una delle
più diffuse, adottata ad esempio in \cite{giusti}, risale a Dedekind (1872) e usa la
nozione di \emph{sezione}, cioè di taglio della retta in due parti. Confrontiamola con
la nostra: il confronto è un buon esercizio sull'uso degli assiomi, e mostra che una
stessa proprietà si può esprimere in modi che a prima vista sembrano diversi.

\begin{definizione}[Sezione di $\R$]\label{def:sezione}\index{sezione!di $\R$}
Una coppia $(A,B)$ di sottoinsiemi di $\R$ si chiama \emph{sezione} di $\R$ se
\begin{enumerate}[label=\textup{(\roman*)}]
  \item $A\neq\emptyset$ e $B\neq\emptyset$;
  \item $A\cup B=\R$ e $A\cap B=\emptyset$;
  \item $a<b$ per ogni $a\in A$ e per ogni $b\in B$.
\end{enumerate}
Un numero $L$ tale che $a\le L\le b$ per ogni $a\in A$ e per ogni $b\in B$ si chiama
\emph{elemento separatore} della sezione.
\end{definizione}

L'\emph{assioma di Dedekind}\index{assioma di Dedekind} afferma:
\begin{center}
\textbf{(D)}\enspace \emph{ogni sezione di $\R$ ha uno e un solo elemento separatore.}
\end{center}
Le due formulazioni non dicono la stessa cosa. Il nostro \cref{ass:completezza} si
applica a una coppia qualsiasi di insiemi non vuoti con $a\le b$: gli insiemi non devono
esaurire $\R$ e possono perfino avere un elemento in comune, come $A=B=\{0\}$; in cambio,
afferma soltanto l'\emph{esistenza} di un separatore. L'assioma (D) si applica solo alle
sezioni, ma ne afferma anche l'\emph{unicità}, che in generale è falsa: per $A=[0,1]$ e
$B=[2,3]$ ogni $c\in[1,2]$ è un separatore. Nessuna delle due formulazioni sembra
contenere l'altra; eppure, in presenza degli altri assiomi, si equivalgono.

\begin{teorema}[Equivalenza con l'assioma di Dedekind]\label{teo:dedekind-equivalente}
Supponiamo che valgano gli assiomi \textup{(A1)--(A6)} e \textup{(O1)--(O4)}. Allora
l'assioma di completezza (\cref{ass:completezza}) e l'assioma \textup{(D)} sono
equivalenti: ciascuno dei due implica l'altro.
\end{teorema}

\begin{schema}
\Ipotesi{Gli assiomi (A1)--(A6) e (O1)--(O4).}
\Tesi{(i) completezza $\implica$ (D); \quad (ii) (D) $\implica$ completezza.}
\end{schema}

\begin{strategia}
Entrambe le implicazioni si reggono su un'unica osservazione: \emph{tra due numeri
distinti $u<v$ ce n'è sempre un terzo}, il punto medio $\frac{u+v}{2}$. (i) Una sezione è
una particolare coppia di insiemi separati, dunque il separatore esiste per il nostro
assioma; se ce ne fossero due, $c<c'$, il loro punto medio non potrebbe stare né in $A$
né in $B$, contro $A\cup B=\R$. (ii) Dati $A$ e $B$ come nel nostro assioma, li
«completiamo» a una sezione: a destra mettiamo \emph{tutti} i numeri che sono $\ge$ di
ogni elemento di $A$, a sinistra tutti gli altri. L'assioma (D) fornisce un separatore
$L$ di questa sezione; resta da controllare che $L$ separi anche $A$ e $B$, e il punto
medio servirà ancora per escludere che qualche elemento di $A$ superi $L$.
\end{strategia}

\begin{dimostrazione}
Premettiamo il fatto annunciato: \emph{se $u<v$ e $m=\frac{u+v}{2}$, allora $u<m<v$}.
Infatti $2=1+1\ge0+1=1>0$ (per \ref{ass:ordine-somma} e l'\cref{es:zero-minore-uno}),
quindi $\frac12>0$ (\cref{es:inverso-positivo}); inoltre $m-u=v-m=\frac{v-u}{2}$ è il
prodotto dei numeri positivi $v-u$ e $\frac12$, che è $\ge0$ per
\ref{ass:ordine-prodotto} e diverso da zero per la \cref{prop:annullamento}. Dunque
$m-u>0$ e $v-m>0$.

(i) Supponiamo che valga l'assioma di completezza e sia $(A,B)$ una sezione di $\R$. Gli
insiemi $A$ e $B$ sono non vuoti e $a<b$, quindi $a\le b$, per ogni $a\in A$ e $b\in B$:
per il \cref{ass:completezza} esiste un elemento separatore $c$. Se $c'$ fosse un secondo
separatore diverso da $c$, potremmo supporre (dicotomia) $c<c'$, e il punto medio $m$ di
$c$ e $c'$ verificherebbe $c<m<c'$. Poiché $A\cup B=\R$, $m$ sta in $A$ oppure in $B$.
Se $m\in A$, essendo $c$ un separatore, $m\le c$, contro $c<m$; se $m\in B$, essendo $c'$
un separatore, $c'\le m$, contro $m<c'$. Dunque il separatore è unico: vale (D).

(ii) Supponiamo che valga (D) e siano $A$, $B$ non vuoti con $a\le b$ per ogni $a\in A$ e
$b\in B$. Poniamo
\[
  B^{*}=\{x\in\R : a\le x \ \perogni a\in A\}, \qquad A^{*}=\R\setminus B^{*}.
\]
Per la dicotomia, $x\in A^{*}$ significa che esiste $a\in A$ con $x<a$. Verifichiamo che
$(A^{*},B^{*})$ è una sezione. $B^{*}$ contiene $B$, dunque non è vuoto; $A^{*}$ non è
vuoto, perché se $a_{0}\in A$ allora $a_{0}-1<a_{0}$, cioè $a_{0}-1\in A^{*}$. La
proprietà (ii) della \cref{def:sezione} vale per costruzione. Infine, se $x\in A^{*}$ e
$y\in B^{*}$, esiste $a\in A$ con $x<a$, ed è $a\le y$ perché $y\in B^{*}$: dunque $x<y$.

Per (D) la sezione $(A^{*},B^{*})$ ha un elemento separatore $L$, cioè $x\le L\le y$ per
ogni $x\in A^{*}$ e $y\in B^{*}$. In particolare $L\le b$ per ogni $b\in B$, perché
$B\subseteq B^{*}$. Resta da provare che $a\le L$ per ogni $a\in A$. Se fosse $L<a$ per
qualche $a\in A$, il punto medio $m$ di $L$ e $a$ verificherebbe $L<m<a$; ma $m<a$
significa $m\in A^{*}$, e quindi $m\le L$: assurdo. Dunque $a\le L\le b$ per ogni $a\in
A$ e $b\in B$, e vale l'assioma di completezza.
\end{dimostrazione}

\begin{marginfigure}
\centering
\begin{tikzpicture}[scale=0.85, every node/.style={font=\scriptsize}]
  \foreach \r/\col/\tit in {0/blu/{$L=\max A$}, 1/rame/{$L=\min B$}} {
  \begin{scope}[yshift=-\r*1.25cm]
    \draw[blu, very thick] (0,0) -- (2.4,0);
    \draw[rame, very thick] (2.4,0) -- (4.8,0);
    \fill[\col] (2.4,0) circle (2.4pt);
    \node[above] at (2.4,0.1) {\tit};
    \node[blu, below] at (1.2,-0.06) {$A$};
    \node[rame, below] at (3.6,-0.06) {$B$};
  \end{scope}}
  \begin{scope}[yshift=-2.5cm]
    \draw[blu, very thick, dotted] (0,0) -- (2.3,0);
    \draw[rame, very thick, dotted] (2.5,0) -- (4.8,0);
    \draw[rosso, thick] (2.4,0) circle (2.4pt);
    \node[above] at (2.4,0.1) {in $\Q$: nessun separatore};
    \node[rosso, below] at (2.4,-0.06) {$\sqrt2$};
  \end{scope}
\end{tikzpicture}
\caption{Le due forme di una sezione di $\R$: il punto di taglio $L$ sta in $A$ (in
alto) oppure in $B$ (al centro). In basso, la sezione di $\Q$ della
\eqref{eq:AB-radice2}: il taglio cade in un «buco».}
\label{fig:sezioni}
\end{marginfigure}

\begin{esempio}[Come è fatta una sezione di $\R$]\label{es:sezioni-R}
Sia $(A,B)$ una sezione di $\R$ e sia $L$ il suo elemento separatore. Se $x<L$, allora
$x\notin B$, perché ogni elemento di $B$ è $\ge L$, e dunque $x\in A$; allo stesso modo,
se $x>L$, allora $x\in B$. Il numero $L$ appartiene a uno solo dei due insiemi, e quindi
le sezioni di $\R$ sono soltanto di due tipi:
\begin{align*}
  A&=\{x\le L\},\ \ B=\{x>L\}, && \text{e allora } L=\max A;\\
  A&=\{x<L\},\ \ B=\{x\ge L\}, && \text{e allora } L=\min B
\end{align*}
(\cref{fig:sezioni}). Ogni taglio della retta reale avviene \emph{in un punto}. In $\Q$
le cose vanno diversamente: gli insiemi \eqref{eq:AB-radice2} formano una sezione di
$\Q$, come abbiamo visto nella dimostrazione del \cref{teo:Q-non-completo}, ma $A$ non ha
massimo e $B$ non ha minimo: il taglio non passa per alcun numero razionale.
\end{esempio}

\begin{osservazione}[Il separatore come «primo dei maggioranti»]
Nella parte (ii) della dimostrazione, $B^{*}$ è l'insieme dei numeri che sono $\ge$ di
tutti gli elementi di $A$: nel prossimo capitolo li chiameremo \emph{maggioranti} di $A$.
Il separatore $L$ è $\ge$ di ogni elemento di $A$, dunque appartiene a $B^{*}$, e non
supera alcun elemento di $B^{*}$: è il più piccolo dei maggioranti, cioè l'\emph{estremo
superiore} di $A$, lo strumento con cui useremo la completezza in tutto il libro
(\cref{teo:esistenza-sup}).
\end{osservazione}

\subsection{Approfondimento: un modello dei numeri reali}
\label{subsec:modello-reali}

Nel nostro approccio l'esistenza di $\R$ è un postulato. Dedekind, nel 1872, percorse la
strada opposta: partendo dai soli numeri razionali \emph{costruì} un insieme che
soddisfa tutti gli assiomi del \cref{sec:assiomi}. Ne diamo un'idea, seguendo
\cite{giusti}, e dimostriamo per intero la proprietà più importante, la completezza. In
questo sottoparagrafo dimentichiamo dunque $\R$ e supponiamo di conoscere soltanto $\Q$,
con le sue operazioni e il suo ordinamento, che soddisfano (A1)--(A6) e (O1)--(O4).
L'idea guida viene dall'\cref{es:sezioni-R}: un numero reale è individuato dal taglio
che produce sui razionali; allora \emph{definiamo} un numero reale come un taglio di
$\Q$.

\begin{definizione}[Sezione di $\Q$]\label{def:sezione-Q}\index{sezione!di $\Q$}
Una \emph{sezione} di $\Q$ è una coppia $(A,B)$ di sottoinsiemi non vuoti di $\Q$ tali
che $A\cup B=\Q$, $A\cap B=\emptyset$, $a<b$ per ogni $a\in A$ e $b\in B$, e tali che
$A$ \emph{non abbia massimo}. Indichiamo con $\widetilde{\R}$ l'insieme di tutte le
sezioni di $\Q$, e diciamo che $\alpha=(A,B)$ precede $\beta=(C,D)$, e scriviamo
$\alpha\preceq\beta$, se $A\subseteq C$.
\end{definizione}

La richiesta che $A$ non abbia massimo serve a evitare doppioni: il taglio prodotto da
un razionale $r$ si può fare mettendo $r$ a sinistra o a destra, e conveniamo di
metterlo sempre a destra. Si noti che $B=\Q\setminus A$: una sezione è determinata dalla
sua classe inferiore $A$, e per questo l'ordinamento si definisce confrontando le classi
inferiori.

\begin{esempio}[Sezioni razionali e irrazionali]\label{es:sezioni-Q}
Per ogni $r\in\Q$ la coppia
\[
  \alpha_{r}=\bigl(\{q\in\Q : q<r\},\ \{q\in\Q : q\ge r\}\bigr)
\]
è una sezione di $\Q$: la classe inferiore non ha massimo, perché se $q<r$ il punto
medio di $q$ e $r$, che è razionale, sta ancora a sinistra di $r$ ed è maggiore di $q$.
In $\alpha_{r}$ il numero $r$ è il minimo della classe superiore. Invece la coppia
$(A,B)$ della \eqref{eq:AB-radice2} è una sezione di $\Q$ in cui $A$ non ha massimo e $B$
non ha minimo (lo abbiamo visto nella dimostrazione del \cref{teo:Q-non-completo}):
non è della forma $\alpha_{r}$ per alcun razionale $r$. Nel modello, questa sezione è il
numero «nuovo» $\sqrt2$.
\end{esempio}

\begin{proposizione}[L'ordinamento del modello]\label{prop:ordine-modello}
La relazione $\preceq$ è una relazione d'ordine totale in $\widetilde{\R}$.
\end{proposizione}

\begin{schema}
\Ipotesi{$\alpha=(A,B)$, $\beta=(C,D)$, $\gamma=(E,F)$ elementi di $\widetilde{\R}$.}
\Tesi{$\alpha\preceq\alpha$; \ se $\alpha\preceq\beta$ e $\beta\preceq\alpha$, allora
$\alpha=\beta$; \ se $\alpha\preceq\beta$ e $\beta\preceq\gamma$, allora
$\alpha\preceq\gamma$; \ $\alpha\preceq\beta$ oppure $\beta\preceq\alpha$.}
\end{schema}

\begin{strategia}
Le proprietà riflessiva, antisimmetrica e transitiva sono quelle dell'inclusione tra
insiemi. L'unica che richiede un'idea è la dicotomia: due sottoinsiemi di $\Q$ in
generale non sono confrontabili per inclusione (si pensi a $\{1\}$ e $\{2\}$). Le classi
inferiori, però, sono «chiuse verso sinistra»: con ogni loro elemento contengono tutti i
razionali più piccoli. Se allora $A\not\subseteq C$ e $a\in A\setminus C$, l'intera
classe $C$ sta a sinistra di $a$, e quindi dentro $A$.
\end{strategia}

\begin{dimostrazione}
La riflessività e la transitività seguono da quelle dell'inclusione. Se $A\subseteq C$ e
$C\subseteq A$, allora $A=C$ e quindi anche $B=\Q\setminus A=\Q\setminus C=D$: dunque
$\alpha=\beta$.

Proviamo preliminarmente che, se $a\in A$ e $q\in\Q$ con $q<a$, allora $q\in A$.
Infatti, se fosse $q\in B$, sarebbe $a<q$ per la definizione di sezione, contro $q<a$; e
poiché $A\cup B=\Q$, deve essere $q\in A$.

Per la dicotomia, supponiamo che non sia $A\subseteq C$ e sia $a\in A$ con $a\notin C$,
cioè $a\in D$. Per ogni $c\in C$ si ha $c<a$, perché $c\in C$ e $a\in D$; per quanto
appena provato, $c\in A$. Dunque $C\subseteq A$, cioè $\beta\preceq\alpha$.
\end{dimostrazione}

\begin{teorema}[Completezza del modello]\label{teo:completezza-modello}
Siano $X$ e $Y$ due sottoinsiemi non vuoti di $\widetilde{\R}$ tali che
$\alpha\preceq\beta$ per ogni $\alpha\in X$ e per ogni $\beta\in Y$. Allora esiste
$\lambda\in\widetilde{\R}$ tale che $\alpha\preceq\lambda\preceq\beta$ per ogni
$\alpha\in X$ e per ogni $\beta\in Y$.
\end{teorema}

\begin{schema}
\Ipotesi{$X,Y\subseteq\widetilde{\R}$ non vuoti; \ $\alpha\preceq\beta$ per ogni
$\alpha\in X$, $\beta\in Y$.}
\Tesi{Esiste $\lambda\in\widetilde{\R}$ con $\alpha\preceq\lambda\preceq\beta$ per ogni
$\alpha\in X$, $\beta\in Y$.}
\end{schema}

\begin{strategia}
Poiché l'ordine è l'inclusione delle classi inferiori, il candidato naturale per la classe
inferiore di $\lambda$ è l'\emph{unione} $A_{0}$ delle classi inferiori degli elementi
di $X$: contiene ciascuna di esse, quindi $\lambda$ seguirà ogni $\alpha\in X$; ed è
contenuta nella classe inferiore di ogni $\beta\in Y$, perché questa contiene già tutte
le classi che stiamo unendo. Bisogna soltanto verificare che $A_{0}$ e il suo
complementare formino una sezione di $\Q$: ciascuna proprietà di $(A_{0},\Q\setminus
A_{0})$ discende dalla proprietà corrispondente delle sezioni di $X$ e di $Y$.
\end{strategia}

\begin{dimostrazione}
Poniamo
\[
  A_{0}=\bigcup_{(A,B)\in X} A, \qquad B_{0}=\Q\setminus A_{0},
\]
cioè $q\in A_{0}$ se e solo se $q\in A$ per almeno una sezione $(A,B)\in X$. Fissiamo
$\beta=(C,D)\in Y$: per ipotesi $A\subseteq C$ per ogni $(A,B)\in X$, e quindi
$A_{0}\subseteq C$. Verifichiamo che $\lambda=(A_{0},B_{0})$ è una sezione di $\Q$.
\begin{itemize}
  \item $A_{0}\neq\emptyset$, perché $X$ non è vuoto e le classi inferiori non sono vuote;
  $B_{0}\neq\emptyset$, perché $A_{0}\subseteq C$ implica $D=\Q\setminus C\subseteq
  B_{0}$, e $D\neq\emptyset$.
  \item $A_{0}\cup B_{0}=\Q$ e $A_{0}\cap B_{0}=\emptyset$ per costruzione.
  \item Siano $a\in A_{0}$ e $b\in B_{0}$. Esiste $(A,B)\in X$ con $a\in A$; poiché
  $b\notin A_{0}$, è $b\notin A$, cioè $b\in B$, e dunque $a<b$.
  \item $A_{0}$ non ha massimo: se $a\in A_{0}$, allora $a\in A$ per qualche $(A,B)\in
  X$; $A$ non ha massimo, dunque esiste $a'\in A$ con $a'>a$, e $a'\in A_{0}$.
\end{itemize}
Infine, $A\subseteq A_{0}$ per ogni $(A,B)\in X$, cioè $\alpha\preceq\lambda$ per ogni
$\alpha\in X$; e $A_{0}\subseteq C$ per ogni $(C,D)\in Y$, cioè $\lambda\preceq\beta$ per
ogni $\beta\in Y$.
\end{dimostrazione}

\begin{osservazione}[Le operazioni nel modello]
Per completare la costruzione bisogna definire in $\widetilde{\R}$ la somma e il prodotto
e verificare gli assiomi (A1)--(A6) e (O1)--(O4). La somma si definisce in modo naturale:
se $\alpha=(A,B)$ e $\beta=(C,D)$, si pone $\alpha+\beta=(A+C,\ \Q\setminus(A+C))$, dove
$A+C=\{a+c : a\in A,\ c\in C\}$, e lo zero è $\alpha_{0}$. Il prodotto è più delicato a
causa della regola dei segni: per due sezioni «positive» si moltiplicano gli elementi
delle classi superiori, non quelli delle classi inferiori, che contengono numeri negativi
grandi in valore assoluto, con prodotti positivi e grandissimi. Le verifiche sono
elementari ma lunghe, e le omettiamo. Si prova infine che
$\alpha_{r}+\alpha_{s}=\alpha_{r+s}$, che $\alpha_{r}\cdot\alpha_{s}=\alpha_{rs}$ e che
$\alpha_{r}\preceq\alpha_{s}$ se e solo se $r\le s$: la corrispondenza $r\mapsto\alpha_{r}$
«immerge» $\Q$ nel modello senza alterarne operazioni e ordinamento, e la sezione
dell'\cref{es:sezioni-Q} risulta avere quadrato $\alpha_{2}$.
\end{osservazione}

\begin{osservazione}[Costruire o postulare?]
Che cosa si guadagna costruendo $\widetilde{\R}$? Weierstrass e i matematici della sua
generazione giudicavano la costruzione più soddisfacente del metodo assiomatico, perché
riconduce l'esistenza dei reali a quella, ritenuta più sicura, dei razionali e quindi
dei naturali (si veda la nota storica del \cref{sec:assiomi}). Ma in matematica
«esistere» significa, in ultima analisi, «non essere contraddittorio», e la costruzione
fa largo uso della teoria degli insiemi: ogni sezione è una coppia di insiemi infiniti di
razionali, e $\widetilde{\R}$ è un insieme di tali coppie. Ciò che si ottiene è dunque una
garanzia \emph{relativa}: se non sono contraddittori gli assiomi dei naturali e quelli
della teoria degli insiemi, non lo sono neppure gli assiomi di $\R$. Vale anche il
viceversa, perché dentro $\R$ c'è un insieme, $\N$, che soddisfa gli assiomi di Peano
(\cref{oss:peano}). Una garanzia assoluta non è raggiungibile: nel 1931 Kurt Gödel
(1906--1978) dimostrò che la non contraddittorietà dell'aritmetica non si può provare con
i soli mezzi dell'aritmetica stessa. La scelta del punto di partenza è quindi, in larga
misura, una questione di convenienza; per l'Analisi, che studia il continuo, è naturale
partire dagli assiomi dei numeri reali \cite{giusti}.
\end{osservazione}
